\documentclass[../../../clio.tex]{subfiles}
\begin{document}

\chapter{Constraint Propagation}
\label{ch:chapter-05-constraint-propagation}

\section{Chapter Preface}
Local decisions can force global consequences. Crossword fills, Sudoku constraints, molecular interactions, and distributed consistency protocols all demonstrate that seemingly isolated commitments reverberate through a system.

Constraint propagation is the machinery that carries those reverberations. Rather than treating propagation as a heuristic, CLIO formulates it as an operator-theoretic process over admissible candidate sets.

This chapter builds the fixed-point perspective needed to reason about termination, convergence, and correctness at scale.

\section{Derivations and Formal Results}
Let each constraint define a pruning operator \(T_i:\mathcal P(\stateSpace)\to\mathcal P(\stateSpace)\), and define
\[
T=T_k\circ\cdots\circ T_1.
\]

\begin{proposition}[Monotone set contraction]
If each \(T_i(S)\subseteq S\), then sequence \(S_{n+1}=T(S_n)\) is monotonically decreasing under set inclusion.
\end{proposition}

\begin{proof}
By composition, \(T(S_n)\subseteq S_n\) for each \(n\), so \(S_{n+1}\subseteq S_n\).
\end{proof}

\begin{proposition}[Finite termination on finite state spaces]
If \(\stateSpace\) is finite, the monotone sequence \(S_{n+1}=T(S_n)\) stabilizes in finitely many steps.
\end{proposition}

\begin{theorem}[Knaster--Tarski fixed points]
If \(T\) is monotone on a complete lattice, then \(T\) has least and greatest fixed points.
\end{theorem}

\begin{theorem}[Banach convergence for metric propagators]
If \((M,d)\) is complete and \(d(Tx,Ty)\le qd(x,y)\) with \(0<q<1\), then \(T\) has a unique fixed point and Picard iteration converges to it.
\end{theorem}

\end{document}
